\section{Runtime model}
\label{sec:model}

A system consists of a finite set of agents $A$, a shared memory $M$ mapping
cells to values with a distinguished sentinel $\mathrm{NULL}$, and a shared
registry $R$ of available tools. An agent executes a sequence of
\emph{operations}, each in three phases: a \emph{read} phase that captures the
values of a chosen subset of cells (the operation's $\rs$) together with the
registry contents ($\rreg$), stamping the read with the current log length as
$\rt$; a \emph{generation} phase, dominated by inference latency and not
treated as instantaneous; and a \emph{write} phase that atomically appends a
record to the log, installing the planned writes ($\ws$, $\wv$), the selected
tool ($\pt$) and the registry at commit ($\wreg$), stamped $\wt$. The log
length is a logical clock~\cite{Lamport1978TimeClocks}; a history $h$ is the
log, $h\in\mathrm{Seq}(\fld{OpRecord})$. The specification is
\texttt{Memory.tla}.

\smallskip\noindent\emph{Store axiom.} A read returns the latest committed
value at or before its read time. TLC confirms this and the two hypotheses
below as invariants of \texttt{Memory.tla} over all
\CommitOrderMemoryStates{} reachable states at the checked bounds
(\cref{sec:a1}).

\smallskip\noindent\emph{Per-agent non-overlap; read before commit.} Two
operations of one agent do not overlap in $[\rt,\wt]$; every operation reads
strictly before it commits. Both are invariants of the model, and both are
load-bearing for \cref{thm:corc} and for the detector transcription of
\cref{sec:verus-detectors}. A framework that runs one node as several parallel
instances under one name (LangGraph's \texttt{Send} fan-out) violates the first
and is outside the model; \cref{sec:prevalence} counts how often that shape
appears in public code.

\smallskip\noindent\emph{Determinism is not assumed by the theorems.} Every
mechanized result ranges over recorded histories with write values free. Under
sampling, a value mismatch need not change the downstream decision, so the
\A1 detector is a sound candidate screen that over-reports by a factor
\cref{sec:materiality} measures. Durable-execution engines take the same
stance~\cite{BurckhardtEtAl2021Durable}: replay is deterministic because every
non-deterministic step's result is recorded, not because the step is.

\smallskip\noindent\emph{Atomic commit is at the record level.} External
effects are outside the store and are governed separately: their order by
\A6 and $L_3$, their release relative to a still-retractable basis by \A3 and
$L_2$. At-most-once externalization is assumed and not modeled; \cref{sec:limitations}
reports a deployed resumption path that violates it.

\smallskip\noindent\emph{No replication.} The memory is one logical store.
Split view under replication (\A4) is formalized separately at the Verus
model level (supplement \S\,J) and plays no role in the lattice.

\smallskip\noindent\emph{Placement.} The model does not derive from any one
runtime. LangGraph is driven through it live (\cref{sec:live-a1}); three
others are placed informally (supplement \S\,D).
